% !TEX TS-program = xelatex
% !TEX encoding = UTF-8 Unicode
% -*- coding: UTF-8; -*-
% vim: set fenc=utf-8

\documentclass[11pt,a4paper]{article}

% XeLaTeX font setup for Vietnamese
\usepackage{fontspec}
\setmainfont{Times New Roman}
\setsansfont{Arial}
\setmonofont{Consolas}

\usepackage[margin=0.55in]{geometry}
\usepackage{hyperref}
\usepackage{xcolor}
\usepackage{enumitem}
\usepackage{titlesec}
\usepackage{fontawesome5}

% Define colors
\definecolor{darkblue}{rgb}{0.0, 0.0, 0.55}
\definecolor{darkgray}{rgb}{0.3, 0.3, 0.3}
\definecolor{accentblue}{HTML}{1A73E8}

% Remove page number
\pagestyle{empty}

% Hyperlink setup
\hypersetup{
  colorlinks=true,
  linkcolor=accentblue,
  urlcolor=accentblue
}

% Section formatting
\titleformat{\section}
  {\large\bfseries\color{darkblue}}
  {}
  {0em}
  {\MakeUppercase}
  [\vspace{-0.6em}\rule{\textwidth}{0.5pt}]

\titlespacing{\section}{0pt}{0.6em}{0.3em}

% Custom commands
\newcommand{\cvname}[1]{
  \centerline{\LARGE\textbf{#1}}
  \vspace{0.3em}
}

\newcommand{\cvcontact}[1]{
  \centerline{\small #1}
}

\newcommand{\cventry}[4]{
  \noindent\textbf{#1} \hfill \textit{#2} \\
  \textit{#3} \\
  #4
  \vspace{0.4em}
}

\begin{document}

% ============================================================
% HEADER
% ============================================================
\cvname{Lê Đức Nguyên}

\cvcontact{
  \faEnvelope\ \href{mailto:nguyen.ld2310@gmail.com}{nguyen.ld2310@gmail.com}
  \quad \faPhone\ +84 945 032 669
  \quad \faGithub\ \href{https://github.com/AndrewHaward2310}{AndrewHaward2310}
  \quad \faMapMarker*\ Hà Nội, Việt Nam
}

\vspace{0.5em}

% ============================================================
% MỤC TIÊU NGHỀ NGHIỆP
% ============================================================
\section{Mục tiêu nghề nghiệp}

Sinh viên năm cuối ngành Cơ điện tử, Đại học Bách khoa Hà Nội, mong muốn ứng tuyển vị trí \textbf{Thực tập sinh} trong lĩnh vực \textbf{Robotics, AI và Hệ thống tự động} tại Vingroup. Có kinh nghiệm thực tế với ROS, Reinforcement Learning, điều khiển PID và hệ thống IoT công nghiệp. Sẵn sàng đóng góp và học hỏi trong môi trường nghiên cứu \& phát triển sản phẩm công nghệ cao.

% ============================================================
% HỌC VẤN
% ============================================================
\section{Học vấn}

\cventry
  {Đại học Bách khoa Hà Nội (HUST)}
  {10/2021 -- 03/2026}
  {Cử nhân Kỹ thuật Cơ điện tử}
  {
    \begin{itemize}[leftmargin=1.5em, topsep=0pt, itemsep=0.15em]
      \item Chuyên ngành: Cơ điện tử, hệ thống tự động hóa, điều khiển nhúng
      \item Các môn liên quan: Điều khiển Tự động, Robotics, Thiết kế Hệ thống Số, Xử lý Tín hiệu, Hệ thống Nhúng
      \item Đạt Học bổng Khuyến khích Học tập (Kỳ 20211, 20212)
    \end{itemize}
  }

% ============================================================
% KINH NGHIỆM DỰ ÁN & NGHIÊN CỨU
% ============================================================
\section{Kinh nghiệm Dự án \& Nghiên cứu}

% --- Project 1: AGV (Ưu tiên cao nhất) ---
\cventry
  {\href{https://andrewhaward2310.github.io/gnn-il-ppo-paper/}{Hệ thống Điều phối Đa Robot AGV bằng Reinforcement Learning}}
  {07/2025 -- 09/2025}
  {Trưởng nhóm / Lập trình viên chính | Viện Cơ khí, ĐHBK Hà Nội}
  {
    \begin{itemize}[leftmargin=1.5em, topsep=0pt, itemsep=0.15em]
      \item Thiết kế và triển khai framework Multi-Agent Reinforcement Learning để điều phối xe tự hành (AGV) trong môi trường nhà máy, xử lý đồng thời nhiều robot với phân công nhiệm vụ tự động
      \item Xây dựng hệ thống điều khiển PID tích hợp RL agent, đạt độ chính xác cao trong điều khiển chuyển động robot
      \item Huấn luyện mô hình RL sử dụng thuật toán PPO (Proximal Policy Optimization) cho bài toán phân bổ nhiệm vụ và tránh va chạm
      \item Sử dụng ROS (Robot Operating System) cho giao tiếp robot và mô phỏng thời gian thực
      \item Phát triển simulation 2D có giao diện web với benchmark so sánh hiệu năng thuật toán
      \item \textbf{Quy mô:} 2 thành viên \quad | \quad \textbf{Công nghệ:} PPO, RL, Imitation Learning, PID Control, ROS, Python, GNN
    \end{itemize}
  }

% --- Project 2: DENSO Factory Hacks ---
\cventry
  {Hệ thống Cảnh báo Lỗi Công nghiệp dựa trên LLM}
  {11/2023 -- 01/2024}
  {Trưởng nhóm | DENSO Factory Hacks 2023 (Hợp tác FPT \& DENSO)}
  {
    \begin{itemize}[leftmargin=1.5em, topsep=0pt, itemsep=0.15em]
      \item Dẫn dắt nhóm 4 thành viên nghiên cứu và triển khai giải pháp phát hiện lỗi công nghiệp sử dụng Large Language Models
      \item Thiết kế kiến trúc hệ thống tích hợp Elasticsearch (indexing dữ liệu) và Weaviate (tìm kiếm ngữ nghĩa - vector database)
      \item Phát triển dashboard thời gian thực bằng Streamlit để giám sát và phân tích lỗi
      \item Cộng tác trực tiếp với mentors từ FPT và DENSO để triển khai giải pháp cấp sản xuất
      \item \textbf{Quy mô:} 4 thành viên \quad | \quad \textbf{Công nghệ:} LLM, Elasticsearch, Weaviate, Docker, Streamlit, Python
      \item \textbf{Giải thưởng:} Giải Ba cuộc thi DENSO Factory Hacks 2023
    \end{itemize}
  }

% --- Project 3: IoT Smart Garden ---
\cventry
  {\href{https://github.com/AndrewHaward2310/IoTSmartGardenSystem}{Hệ thống Vườn Thông minh IoT}}
  {10/2025 -- 01/2026}
  {Trưởng nhóm Backend / Kiến trúc sư Hệ thống | Đồ án Môn học}
  {
    \begin{itemize}[leftmargin=1.5em, topsep=0pt, itemsep=0.15em]
      \item Thiết kế toàn bộ kiến trúc backend cho hệ thống IoT vườn thông minh, từ thu thập dữ liệu đến API phục vụ người dùng
      \item Xây dựng hệ thống microservice có khả năng mở rộng, đóng gói từng service bằng Docker
      \item Triển khai pipeline dữ liệu tích hợp PostgreSQL, MongoDB (time-series sensor data) và Redis (caching hiệu năng cao)
      \item Phát triển MQTT handler service xử lý luồng dữ liệu thời gian thực từ thiết bị ESP32
      \item \textbf{Quy mô:} 4 thành viên \quad | \quad \textbf{Công nghệ:} Node.js, Express.js, PostgreSQL, MongoDB, Redis, Docker, Docker Compose, MQTT, ESP32
    \end{itemize}
  }

% --- Project 4: UWB Warehouse ---
\cventry
  {Hệ thống Quản lý Kho hàng sử dụng Công nghệ UWB}
  {11/2022 -- 03/2023}
  {Lập trình viên | Giải thưởng Nghiên cứu Khoa học \& Sáng tạo}
  {
    \begin{itemize}[leftmargin=1.5em, topsep=0pt, itemsep=0.15em]
      \item Xây dựng kiến trúc database MongoDB cho lưu trữ và quản lý dữ liệu cảm biến UWB
      \item Xử lý dữ liệu thời gian thực từ cảm biến UWB để theo dõi vị trí và quản lý tài sản trong kho
      \item Tích hợp IoT sensors với cloud backend qua giao thức MQTT
      \item \textbf{Quy mô:} 4 thành viên \quad | \quad \textbf{Công nghệ:} MongoDB, FastAPI, Pandas, MQTT, UWB Sensors, Python
    \end{itemize}
  }

% ============================================================
% GIẢI THƯỞNG & HỌC BỔNG
% ============================================================
\section{Giải thưởng \& Học bổng}

\noindent\textbf{Học bổng Khuyến khích Học tập (Dean's List)} \hfill \textit{2021 -- 2022} \\
Viện Cơ khí, Đại học Bách khoa Hà Nội — Đạt thành tích học tập xuất sắc kỳ 20211 và 20212

\vspace{0.3em}

\noindent\textbf{Giải thưởng Nghiên cứu Khoa học \& Sáng tạo} \hfill \textit{06/2022} \\
Đại học Bách khoa Hà Nội — Nghiên cứu sáng tạo trong thiết kế và triển khai hệ thống kỹ thuật

\vspace{0.3em}

\noindent\textbf{Giải Ba, DENSO Factory Hacks 2023} \hfill \textit{01/2024} \\
Tổ chức bởi DENSO và FPT — Hệ thống phát hiện lỗi công nghiệp dựa trên LLM

% ============================================================
% KỸ NĂNG
% ============================================================
\section{Kỹ năng}

\noindent\textbf{Lập trình:} Python (Nâng cao), C/C++ (Trung cấp), MATLAB (Nâng cao), JavaScript/Node.js

\vspace{0.15em}
\noindent\textbf{Robotics \& Điều khiển:} ROS (Robot Operating System), Điều khiển PID, Giao thức MQTT, UWB, Hệ thống nhúng (ESP32), IoT công nghiệp

\vspace{0.15em}
\noindent\textbf{AI \& Machine Learning:} Reinforcement Learning (PPO, Imitation Learning), Graph Neural Networks, Large Language Models, Elasticsearch, Weaviate (Vector DB)

\vspace{0.15em}
\noindent\textbf{DevOps \& Hệ thống:} Docker, Docker Compose, Git, Kiến trúc Microservices, Thiết kế Hệ thống

\vspace{0.15em}
\noindent\textbf{Cơ sở dữ liệu:} MongoDB, PostgreSQL, Redis, Pandas, ETL Pipeline, Xử lý dữ liệu thời gian thực

% ============================================================
% NGƯỜI THAM CHIẾU
% ============================================================
\section{Người tham chiếu}

\noindent\textbf{PGS. TS. Phạm Đức An} \\
Giảng viên Cao cấp, Viện Cơ khí \\
Phó Trưởng Bộ môn Cơ điện tử \\
Đại học Bách khoa Hà Nội (HUST) \\
\href{mailto:an.phamduc@hust.edu.vn}{an.phamduc@hust.edu.vn}

\end{document}
